% GENERATED FILE. Edit canonical lessons, then run npm run generate:books.
% canonical-source-hash: fnv1a64:af5e011a9bc15704
% canonical-lessons: HI-C335-raksas, HI-C335-devta, HI-C335-vardan, HI-C335-amar, HI-C335-nigalna, HI-R335-review-328-335

\chapter{Gods and Demons}
\label{ch:hi-gods-and-demons}
\hlchaptermodality{335}

\begin{chapteropening}
\textbf{By the end of this chapter:} \emph{I can tell a story of gods and demons.}

Complete the last lesson of chapter 335: I can tell a story of gods and demons.
\end{chapteropening}

\section[\texorpdfstring{\hi{राक्षस} (rākṣas)}{राक्षस (rākṣas)}]{\hi{राक्षस} (rākṣas) — a demon}

\label{lesson:HI-C335-raksas}



\begin{quote}
Before the new one: say the Hindi for a riddle, then the Hindi for a saying.
\end{quote}

\subsection*{You'll want to know: \hi{राक्षस}}
\textbf{\hi{राक्षस}} — \emph{rākṣas} — \textquotedblleft{}a demon\textquotedblright{}.

\begin{quote}
\hi{एक} \hi{राक्षस} \hi{जंगल} \hi{में} \hi{रहता} \hi{था।}
\end{quote}

\begin{quote}
\emph{ek rākṣas jaṅgal meṁ rahtā thā.}
\end{quote}

\begin{quote}
A demon lived in the forest.
\end{quote}

Masculine.

\begin{cousinweb}
\textbf{\hi{राक्षस}} (\emph{rākṣas}) is Sanskrit \emph{rākṣasa}, \textquotedblleft{}a demon\textquotedblright{} (Monier-Williams): the giant, man-eating enemy of the old tales.
\end{cousinweb}

\begin{grammarlens}[title={What you've built}]
One more word for telling what happened.
\end{grammarlens}

\subsection*{Your turn}
\begin{itemize}
\raggedright
  \item \emph{Say it:} \emph{rākṣas}
  \item \emph{Say it:} \emph{rākṣas}, once more
  \item \emph{Say it:} \emph{ek rākṣas jaṅgal meṁ rahtā thā.} — A demon lived in the forest.
  \item \emph{Recall:} say the Hindi for a riddle, then the Hindi for a saying, then say \emph{rākṣas} again
\end{itemize}

\begin{tcolorbox}[breakable,colback=teal!4,colframe=teal!35!black,title={Before you move on}]
What does \hi{राक्षस} mean? (\textquotedblleft{}A demon\textquotedblright{}.) Say it once more.
\end{tcolorbox}
\section[\texorpdfstring{\hi{देवता} (devtā)}{देवता (devtā)}]{\hi{देवता} (devtā) — a god; a deity}

\label{lesson:HI-C335-devta}



\begin{quote}
Before the new one: say the Hindi for a saying, then the Hindi for a demon.
\end{quote}

\subsection*{You'll want to know: \hi{देवता}}
\textbf{\hi{देवता}} — \emph{devtā} — \textquotedblleft{}a god; a deity\textquotedblright{}.

\begin{quote}
\hi{कथा} \hi{में} \hi{देवता} \hi{और} \hi{राक्षस} \hi{लड़े।}
\end{quote}

\begin{quote}
\emph{kathā meṁ devtā aur rākṣas laṛe.}
\end{quote}

\begin{quote}
In the tale the gods and the demons fought.
\end{quote}

Masculine.

\begin{cousinweb}
\textbf{\hi{देवता}} (\emph{devtā}) is Sanskrit \emph{devatā}, \textquotedblleft{}divinity, a deity\textquotedblright{} (Monier-Williams). In Hindi it is masculine and means a god.
\end{cousinweb}

\begin{grammarlens}[title={What you've built}]
One more word for telling what happened.
\end{grammarlens}

\subsection*{Your turn}
\begin{itemize}
\raggedright
  \item \emph{Say it:} \emph{devtā}
  \item \emph{Say it:} \emph{devtā}, once more
  \item \emph{Say it:} \emph{kathā meṁ devtā aur rākṣas laṛe.} — In the tale the gods and the demons fought.
  \item \emph{Recall:} say the Hindi for a saying, then the Hindi for a demon, then say \emph{devtā} again
\end{itemize}

\begin{tcolorbox}[breakable,colback=teal!4,colframe=teal!35!black,title={Before you move on}]
What does \hi{देवता} mean? (\textquotedblleft{}A god; a deity\textquotedblright{}.) Say it once more.
\end{tcolorbox}
\section[\texorpdfstring{\hi{वरदान} (vardān)}{वरदान (vardān)}]{\hi{वरदान} (vardān) — a boon}

\label{lesson:HI-C335-vardan}



\begin{quote}
Before the new one: say the Hindi for a demon, then the Hindi for a god.
\end{quote}

\subsection*{You'll want to know: \hi{वरदान}}
\textbf{\hi{वरदान}} — \emph{vardān} — \textquotedblleft{}a boon\textquotedblright{}.

\begin{quote}
\hi{देवता} \hi{ने} \hi{राक्षस} \hi{को} \hi{एक} \hi{वरदान} \hi{दिया।}
\end{quote}

\begin{quote}
\emph{devtā ne rākṣas ko ek vardān diyā.}
\end{quote}

\begin{quote}
The god granted the demon a boon.
\end{quote}

Masculine.

\begin{cousinweb}
\textbf{\hi{वरदान}} (\emph{vardān}) is Sanskrit \emph{varadāna} (Monier-Williams), \emph{vara}, a wish or a choice, and \emph{dāna}, a giving: a wish granted by a god.
\end{cousinweb}

\begin{grammarlens}[title={What you've built}]
One more word for telling what happened.
\end{grammarlens}

\subsection*{Your turn}
\begin{itemize}
\raggedright
  \item \emph{Say it:} \emph{vardān}
  \item \emph{Say it:} \emph{vardān}, once more
  \item \emph{Say it:} \emph{devtā ne rākṣas ko ek vardān diyā.} — The god granted the demon a boon.
  \item \emph{Recall:} say the Hindi for a demon, then the Hindi for a god, then say \emph{vardān} again
\end{itemize}

\begin{tcolorbox}[breakable,colback=teal!4,colframe=teal!35!black,title={Before you move on}]
What does \hi{वरदान} mean? (\textquotedblleft{}A boon\textquotedblright{}.) Say it once more.
\end{tcolorbox}
\section[\texorpdfstring{\hi{अमर} (amar)}{अमर (amar)}]{\hi{अमर} (amar) — immortal}

\label{lesson:HI-C335-amar}



\begin{quote}
Before the new one: say the Hindi for a god, then the Hindi for a boon.
\end{quote}

\subsection*{You'll want to know: \hi{अमर}}
\textbf{\hi{अमर}} — \emph{amar} — \textquotedblleft{}immortal\textquotedblright{}.

\begin{quote}
\hi{वरदान} \hi{से} \hi{राक्षस} \hi{अमर} \hi{हो} \hi{गया।}
\end{quote}

\begin{quote}
\emph{vardān se rākṣas amar ho gayā.}
\end{quote}

\begin{quote}
The boon made the demon immortal.
\end{quote}

\begin{cousinweb}
\textbf{\hi{अमर}} (\emph{amar}) is Sanskrit \emph{amara}, \textquotedblleft{}undying\textquotedblright{} (Monier-Williams), \emph{a-}, not, and the family of \emph{mṛ}, to die, the root of \textbf{\hi{मरना}}.
\end{cousinweb}

\begin{grammarlens}[title={What you've built}]
One more word for telling what happened.
\end{grammarlens}

\subsection*{Your turn}
\begin{itemize}
\raggedright
  \item \emph{Say it:} \emph{amar}
  \item \emph{Say it:} \emph{amar}, once more
  \item \emph{Say it:} \emph{vardān se rākṣas amar ho gayā.} — The boon made the demon immortal.
  \item \emph{Recall:} say the Hindi for a god, then the Hindi for a boon, then say \emph{amar} again
\end{itemize}

\begin{tcolorbox}[breakable,colback=teal!4,colframe=teal!35!black,title={Before you move on}]
What does \hi{अमर} mean? (\textquotedblleft{}Immortal\textquotedblright{}.) Say it once more.
\end{tcolorbox}
\section[\texorpdfstring{\hi{निगलना} (nigalnā)}{निगलना (nigalnā)}]{\hi{निगलना} (nigalnā) — to swallow}

\label{lesson:HI-C335-nigalna}



\begin{quote}
Before the new one: say the Hindi for a boon, then the Hindi for immortal.
\end{quote}

\subsection*{You'll want to know: \hi{निगलना}}
\textbf{\hi{निगलना}} — \emph{nigalnā} — \textquotedblleft{}to swallow\textquotedblright{}.

\begin{quote}
\hi{राक्षस} \hi{ने} \hi{पूरी} \hi{नदी} \hi{निगल} \hi{ली।}
\end{quote}

\begin{quote}
\emph{rākṣas ne pūrī nadī nigal lī.}
\end{quote}

\begin{quote}
The demon swallowed the whole river.
\end{quote}

Past, with \hi{ने} (he or she …): \hi{उसने} \hi{फल} \hi{निगला}, \emph{usne phal niglā}.

\begin{cousinweb}
\textbf{\hi{निगलना}} (\emph{nigalnā}) goes back to Sanskrit \emph{nigalati}, \textquotedblleft{}swallows\textquotedblright{} (CDIAL 7163). It is swallowing whole, without chewing: \textbf{\hi{निगल} \hi{लिया}}.
\end{cousinweb}

\begin{grammarlens}[title={What you've built}]
One more word for telling what happened.
\end{grammarlens}

\subsection*{Your turn}
\begin{itemize}
\raggedright
  \item \emph{Say it:} \emph{nigalnā}
  \item \emph{Say it:} \emph{nigalnā}, once more
  \item \emph{Say it:} \emph{rākṣas ne pūrī nadī nigal lī.} — The demon swallowed the whole river.
  \item \emph{Recall:} say the Hindi for a boon, then the Hindi for immortal, then say \emph{nigalnā} again
\end{itemize}

\begin{tcolorbox}[breakable,colback=teal!4,colframe=teal!35!black,title={Before you move on}]
What does \hi{निगलना} mean? (\textquotedblleft{}To swallow\textquotedblright{}.) Say it once more.
\end{tcolorbox}
\section[(dialogue)]{Looking back: from A Night of Lights to Gods and Demons}

\label{lesson:HI-R335-review-328-335}



\begin{quote}
Say the Hindi for a rangoli, a coloured pattern on the floor and to glitter.
\end{quote}

\subsection*{Your turn}
Say these aloud:

\begin{itemize}
\raggedright
  \item a rangoli, a coloured pattern on the floor, to glitter, to twinkle, a firecracker, fireworks
  \item to sprinkle, to drench, to toss up, to sway (with joy or music), merriment
  \item gathered together, a tableau, to swing, a festive procession, a devotee
  \item a fast, an arati, worship with a lamp, a hymn, a blessed offering (of food), a festive dish
  \item a tradition, a custom, a celebration, a bonfire, an effigy
  \item a ghat, steps down to a river, immersion (of an image), to dip, to be scattered, the floating of lamps
  \item a tale, a sage, a hermitage, a riddle, a saying
  \item a demon, a god, a boon, immortal, to swallow
\end{itemize}

\begin{tcolorbox}[breakable,colback=teal!4,colframe=teal!35!black,title={Before you move on}]
A rangoli, a coloured pattern on the floor? (\textbf{\hi{रंगोली}}, \emph{raṅgolī}.) To glitter? (\textbf{\hi{जगमगाना}}, \emph{jagmagānā}.) To twinkle? (\textbf{\hi{झिलमिलाना}}, \emph{jhilmilānā}.) A firecracker? (\textbf{\hi{पटाख़ा}}, \emph{paṭākhā}.) Fireworks? (\textbf{\hi{आतिशबाज़ी}}, \emph{ātiśbāzī}.) To sprinkle? (\textbf{\hi{छिड़कना}}, \emph{chiṛaknā}.) To drench? (\textbf{\hi{भिगोना}}, \emph{bhigonā}.) To toss up? (\textbf{\hi{उछालना}}, \emph{uchālnā}.) To sway (with joy or music)? (\textbf{\hi{झूमना}}, \emph{jhūmnā}.) Merriment? (\textbf{\hi{मस्ती}}, \emph{mastī}.) Gathered together? (\textbf{\hi{इकट्ठा}}, \emph{ikaṭṭhā}.) A tableau? (\textbf{\hi{झाँकी}}, \emph{jhā̃kī}.) To swing? (\textbf{\hi{झूलना}}, \emph{jhūlnā}.) A festive procession? (\textbf{\hi{शोभायात्रा}}, \emph{śobhāyātrā}.) A devotee? (\textbf{\hi{भक्त}}, \emph{bhakt}.) A fast? (\textbf{\hi{व्रत}}, \emph{vrat}.) An arati, worship with a lamp? (\textbf{\hi{आरती}}, \emph{ārtī}.) A hymn? (\textbf{\hi{भजन}}, \emph{bhajan}.) A blessed offering (of food)? (\textbf{\hi{प्रसाद}}, \emph{prasād}.) A festive dish? (\textbf{\hi{पकवान}}, \emph{pakvān}.) A tradition? (\textbf{\hi{परंपरा}}, \emph{paramparā}.) A custom? (\textbf{\hi{रिवाज}}, \emph{rivāj}.) A celebration? (\textbf{\hi{उत्सव}}, \emph{utsav}.) A bonfire? (\textbf{\hi{अलाव}}, \emph{alāv}.) An effigy? (\textbf{\hi{पुतला}}, \emph{putlā}.) A ghat, steps down to a river? (\textbf{\hi{घाट}}, \emph{ghāṭ}.) Immersion (of an image)? (\textbf{\hi{विसर्जन}}, \emph{visarjan}.) To dip? (\textbf{\hi{डुबोना}}, \emph{ḍubonā}.) To be scattered? (\textbf{\hi{बिखरना}}, \emph{bikharnā}.) The floating of lamps? (\textbf{\hi{दीपदान}}, \emph{dīpdān}.) A tale? (\textbf{\hi{कथा}}, \emph{kathā}.) A sage? (\textbf{\hi{ऋषि}}, \emph{ṛṣi}.) A hermitage? (\textbf{\hi{आश्रम}}, \emph{āśram}.) A riddle? (\textbf{\hi{पहेली}}, \emph{pahelī}.) A saying? (\textbf{\hi{कहावत}}, \emph{kahāvat}.) A demon? (\textbf{\hi{राक्षस}}, \emph{rākṣas}.) A god? (\textbf{\hi{देवता}}, \emph{devtā}.) A boon? (\textbf{\hi{वरदान}}, \emph{vardān}.) Immortal? (\textbf{\hi{अमर}}, \emph{amar}.) To swallow? (\textbf{\hi{निगलना}}, \emph{nigalnā}.)
\end{tcolorbox}
